\section{P04 - Cours - Types construits complément}

\subsection{Objectifs spécifiques}

\begin{itemize}
\tightlist
\item
  Objectif O1 - Identifier précisément la représentation ou la structure
  en jeu.
\item
  Objectif O2 - Appliquer une méthode disciplinaire complète.
\item
  Objectif O3 - Justifier le résultat sur un cas différent.
\item
  Objectif O4 - Contrôler un cas limite et corriger une erreur observée.
\end{itemize}

\subsection{Capacités officielles atomiques}

\begin{itemize}
\tightlist
\item
  P-DATA-CONSTR-01 : Écrire une fonction renvoyant un p-uplet de valeurs
\item
  P-DATA-CONSTR-02B : Construire un tableau par compréhension
\item
  P-DATA-CONSTR-02C : Utiliser des tableaux de tableaux pour représenter
  des matrices
\item
  P-DATA-CONSTR-03A : Construire une entrée de dictionnaire
\item
  P-DATA-CONSTR-03B : Itérer sur les éléments d'un dictionnaire
\item
  P-DATA-CONSTR-03C : Utiliser les méthodes keys, values et items
\end{itemize}

\subsection{Prérequis}

\begin{itemize}
\tightlist
\item
  Connaître les types de base en Python (int, float, str, bool).
\item
  Savoir définir et appeler une fonction avec \texttt{def} et
  \texttt{return}.
\item
  Connaître la syntaxe de base des tuples, listes et dictionnaires
  (séquence P04 initiale).
\item
  Maîtriser les boucles \texttt{for} et \texttt{while}.
\end{itemize}

\subsection{Séance(s) correspondante(s)}

\begin{itemize}
\tightlist
\item
  P04-S1 à P04-S7 : support rattaché aux séances prêtes de la
  progression.
\end{itemize}

\begin{center}\rule{0.5\linewidth}{0.5pt}\end{center}

\subsection{1. P-DATA-CONSTR-01 : Écrire une fonction renvoyant un p-uplet de valeurs}

\subsubsection{Définition}

Un \textbf{p-uplet} (tuple) est une collection ordonnée et
\textbf{immuable} de valeurs. Une fonction peut renvoyer plusieurs
valeurs simultanément en les regroupant dans un p-uplet grâce à
l'instruction \texttt{return}.

La syntaxe \texttt{return\ a,\ b} est équivalente à
\texttt{return\ (a,\ b)}. Python construit automatiquement un tuple.

\subsubsection{Formalisation}

\begin{python}
def nom_fonction(parametres):
    # calculs
    return valeur1, valeur2  # renvoie un tuple (valeur1, valeur2)
\end{python}

Le tuple renvoyé peut être déballé lors de l'appel :

\begin{python}
x, y = nom_fonction(arguments)
\end{python}

\subsubsection{Exemple corrigé 1 - Fonction renvoyant min et max}

\begin{python}
def min_max(tableau):
    """Renvoie le minimum et le maximum d’un tableau non vide."""
    mini = tableau[0]
    maxi = tableau[0]
    for val in tableau:
        if val < mini:
            mini = val
        if val > maxi:
            maxi = val
    return mini, maxi
\end{python}

\begin{itemize}
\tightlist
\item
  Donnée étudiée : \texttt{{[}7,\ 2,\ 9,\ 4,\ 1{]}}.
\item
  Méthode : parcours complet pour trouver les deux extrema.
\item
  Résultat obtenu : \texttt{(1,\ 9)}.
\item
  Déballage :
  \texttt{plus\_petit,\ plus\_grand\ =\ min\_max({[}7,\ 2,\ 9,\ 4,\ 1{]})}.
\end{itemize}

\subsubsection{Exemple corrigé 2 - Coordonnées d’un point milieu}

\begin{python}
def milieu(xa, ya, xb, yb):
    """Renvoie les coordonnées du milieu de [AB]."""
    mx = (xa + xb) / 2
    my = (ya + yb) / 2
    return mx, my
\end{python}

\begin{itemize}
\tightlist
\item
  Donnée étudiée : \texttt{A(0,\ 0)} et \texttt{B(6,\ 4)}.
\item
  Résultat obtenu : \texttt{(3.0,\ 2.0)}.
\end{itemize}

\subsubsection{Cas limites}

\begin{itemize}
\tightlist
\item
  Tableau d'un seul élément pour \texttt{min\_max} : renvoie
  \texttt{(e,\ e)} avec le même élément en min et max.
\item
  Tableau vide : provoque une erreur \texttt{IndexError} car
  \texttt{tableau{[}0{]}} est impossible. Il convient de documenter la
  précondition ou de lever une exception explicite.
\end{itemize}

\begin{center}\rule{0.5\linewidth}{0.5pt}\end{center}

\subsection{2. P-DATA-CONSTR-02B : Construire un tableau par compréhension}

\subsubsection{Définition}

Un \textbf{tableau par compréhension} (list comprehension) permet de
construire une liste en une seule expression lisible. On se limite aux
formes simples et lisibles.

\subsubsection{Formalisation}

\begin{python}
nouveau_tableau = [expression for variable in iterable]
\end{python}

Avec filtre optionnel :

\begin{python}
nouveau_tableau = [expression for variable in iterable if condition]
\end{python}

\subsubsection{Exemple corrigé 1 - Carrés des entiers}

\begin{python}
carres = [x ** 2 for x in range(10)]
# Résultat : [0, 1, 4, 9, 16, 25, 36, 49, 64, 81]
\end{python}

\begin{itemize}
\tightlist
\item
  Donnée étudiée : les entiers de 0 à 9.
\item
  Méthode : appliquer l'opération \texttt{x\ **\ 2} à chaque entier.
\item
  Résultat obtenu : liste de 10 carrés parfaits.
\end{itemize}

\subsubsection{Exemple corrigé 2 - Filtrer les valeurs paires}

\begin{python}
pairs = [n for n in range(20) if n % 2 == 0]
# Résultat : [0, 2, 4, 6, 8, 10, 12, 14, 16, 18]
\end{python}

\subsubsection{Exemple corrigé 3 - Conversion de températures}

\begin{python}
celsius = [0, 10, 20, 30, 40]
fahrenheit = [c * 9 / 5 + 32 for c in celsius]
# Résultat : [32.0, 50.0, 68.0, 86.0, 104.0]
\end{python}

\subsubsection{Cas limites}

\begin{itemize}
\tightlist
\item
  Compréhension sur un itérable vide :
  \texttt{{[}x\ for\ x\ in\ {[}{]}{]}} produit \texttt{{[}{]}}.
\item
  Expression constante : \texttt{{[}0\ for\ \_\ in\ range(5){]}} produit
  \texttt{{[}0,\ 0,\ 0,\ 0,\ 0{]}}, utile pour initialiser un tableau.
\end{itemize}

\subsubsection{Remarque importante}

On ne demande pas d'écrire des compréhensions imbriquées complexes. La
lisibilité prime : si l'expression devient difficile à lire, on préfère
une boucle classique.

\begin{center}\rule{0.5\linewidth}{0.5pt}\end{center}

\subsection{3. P-DATA-CONSTR-02C : Utiliser des tableaux de tableaux pour représenter des matrices}

\subsubsection{Définition}

Une \textbf{matrice} est représentée en Python par une \textbf{liste de
listes} (tableau de tableaux). Chaque sous-liste représente une ligne.
L'accès à l'élément en ligne \texttt{i}, colonne \texttt{j} se fait par
\texttt{matrice{[}i{]}{[}j{]}}. On n'utilise pas de bibliothèque externe
(pas de NumPy).

\subsubsection{Formalisation}

Une matrice de \texttt{n} lignes et \texttt{p} colonnes :

\begin{python}
matrice = [
    [a00, a01, a02],  # ligne 0
    [a10, a11, a12],  # ligne 1
    [a20, a21, a22],  # ligne 2
]
\end{python}

\begin{itemize}
\tightlist
\item
  \texttt{matrice{[}i{]}} donne la ligne \texttt{i} (une liste).
\item
  \texttt{matrice{[}i{]}{[}j{]}} donne l'élément en ligne \texttt{i},
  colonne \texttt{j}.
\item
  \texttt{len(matrice)} donne le nombre de lignes.
\item
  \texttt{len(matrice{[}0{]})} donne le nombre de colonnes.
\end{itemize}

\subsubsection{Exemple corrigé 1 - Création et accès}

\begin{python}
notes = [
    [12, 15, 8],
    [14, 10, 16],
    [9, 11, 13],
]
\end{python}

\begin{itemize}
\tightlist
\item
  \texttt{notes{[}0{]}{[}1{]}} vaut \texttt{15} (ligne 0, colonne 1).
\item
  \texttt{notes{[}2{]}{[}2{]}} vaut \texttt{13} (ligne 2, colonne 2).
\end{itemize}

\subsubsection{Exemple corrigé 2 - Parcours complet d’une matrice}

\begin{python}
def somme_matrice(matrice):
    """Renvoie la somme de tous les éléments de la matrice."""
    total = 0
    for i in range(len(matrice)):
        for j in range(len(matrice[i])):
            total = total + matrice[i][j]
    return total
\end{python}

\begin{itemize}
\tightlist
\item
  Donnée étudiée : la matrice \texttt{notes} ci-dessus.
\item
  Résultat obtenu :
  \texttt{12\ +\ 15\ +\ 8\ +\ 14\ +\ 10\ +\ 16\ +\ 9\ +\ 11\ +\ 13\ =\ 108}.
\end{itemize}

\subsubsection{Exemple corrigé 3 - Initialisation par compréhension}

\begin{python}
zeros = [[0 for j in range(4)] for i in range(3)]
# Résultat : [[0, 0, 0, 0], [0, 0, 0, 0], [0, 0, 0, 0]]
\end{python}

\subsubsection{Cas limites et erreur fréquente}

\textbf{Erreur classique} : initialiser avec la multiplication.

\begin{python}
# INCORRECT : les 3 lignes partagent le même objet en mémoire
mauvais = [[0] * 4] * 3
mauvais[0][0] = 1
# mauvais vaut [[1, 0, 0, 0], [1, 0, 0, 0], [1, 0, 0, 0]]
\end{python}

La bonne méthode est d'utiliser la compréhension
\texttt{{[}{[}0{]}*4\ for\ \_\ in\ range(3){]}}.

\begin{itemize}
\tightlist
\item
  Matrice 1x1 : \texttt{{[}{[}5{]}{]}}, accès
  \texttt{matrice{[}0{]}{[}0{]}}.
\item
  Indice hors limites : \texttt{matrice{[}3{]}{[}0{]}} sur une matrice
  3x3 lève \texttt{IndexError}.
\end{itemize}

\begin{center}\rule{0.5\linewidth}{0.5pt}\end{center}

\subsection{4. P-DATA-CONSTR-03A : Construire une entrée de dictionnaire}

\subsubsection{Définition}

Un \textbf{dictionnaire} associe des \textbf{clés} à des
\textbf{valeurs}. Construire une entrée consiste à ajouter ou modifier
une association \texttt{clé:\ valeur} dans un dictionnaire existant,
avec la syntaxe \texttt{d{[}cle{]}\ =\ valeur}.

\subsubsection{Formalisation}

\begin{python}
d = {}              # dictionnaire vide
d[cle] = valeur     # ajoute ou met à jour l'entrée
\end{python}

Les clés doivent être de type immuable (str, int, tuple). Les valeurs
peuvent être de tout type.

\subsubsection{Exemple corrigé 1 - Métadonnées EXIF d’une photo}

\begin{python}
exif = {}
exif["appareil"] = "Nikon D3500"
exif["date"] = "2025-06-15"
exif["resolution"] = (6000, 4000)
exif["iso"] = 400
exif["flash"] = False
\end{python}

\begin{itemize}
\tightlist
\item
  Donnée étudiée : les métadonnées fictives d'un fichier photo.
\item
  Méthode : chaque affectation \texttt{exif{[}cle{]}\ =\ valeur} crée
  une entrée.
\item
  Résultat obtenu :
  \texttt{\{"appareil":\ "Nikon\ D3500",\ "date":\ "2025-06-15",\ "resolution":\ (6000,\ 4000),\ "iso":\ 400,\ "flash":\ False\}}.
\end{itemize}

\subsubsection{Exemple corrigé 2 - Mise à jour d’une entrée existante}

\begin{python}
exif["iso"] = 800  # modification de la valeur associée à la clé "iso"
\end{python}

Si la clé existe déjà, la valeur est remplacée. Il n'y a pas de doublon
de clé dans un dictionnaire.

\subsubsection{Exemple corrigé 3 - Construction à partir de données}

\begin{python}
def compter_occurrences(texte):
    """Construit un dictionnaire des occurrences de chaque caractère."""
    compteur = {}
    for c in texte:
        if c in compteur:
            compteur[c] = compteur[c] + 1
        else:
            compteur[c] = 1
    return compteur
\end{python}

\begin{itemize}
\tightlist
\item
  Donnée étudiée : \texttt{"abracadabra"}.
\item
  Résultat obtenu :
  \texttt{\{\textquotesingle{}a\textquotesingle{}:\ 5,\ \textquotesingle{}b\textquotesingle{}:\ 2,\ \textquotesingle{}r\textquotesingle{}:\ 2,\ \textquotesingle{}c\textquotesingle{}:\ 1,\ \textquotesingle{}d\textquotesingle{}:\ 1\}}.
\end{itemize}

\subsubsection{Cas limites}

\begin{itemize}
\tightlist
\item
  Clé de type mutable (liste) : \texttt{d{[}{[}1,\ 2{]}{]}\ =\ "x"} lève
  \texttt{TypeError:\ unhashable\ type:\ \textquotesingle{}list\textquotesingle{}}.
\item
  Dictionnaire vide : \texttt{len(\{\})} vaut \texttt{0}, aucune clé
  n'existe.
\end{itemize}

\begin{center}\rule{0.5\linewidth}{0.5pt}\end{center}

\subsection{5. P-DATA-CONSTR-03B : Itérer sur les éléments d’un dictionnaire}

\subsubsection{Définition}

Itérer sur un dictionnaire signifie parcourir ses éléments à l'aide
d'une boucle \texttt{for}. Python offre plusieurs méthodes pour choisir
ce que l'on parcourt.

\subsubsection{Formalisation}

\begin{longtable}[]{@{}ll@{}}
\toprule\noalign{}
Syntaxe & Ce qui est parcouru \\
\midrule\noalign{}
\endhead
\bottomrule\noalign{}
\endlastfoot
\texttt{for\ k\ in\ d} & les clés \\
\texttt{for\ k\ in\ d.keys()} & les clés (explicite) \\
\texttt{for\ v\ in\ d.values()} & les valeurs \\
\texttt{for\ k,\ v\ in\ d.items()} & les couples (clé, valeur) \\
\end{longtable}

\subsubsection{Exemple corrigé 1 - Parcours des clés}

\begin{python}
exif = {"appareil": "Nikon D3500", "date": "2025-06-15", "iso": 400}

for cle in exif:
    print(cle)
# Affiche : appareil, date, iso (un par ligne)
\end{python}

\subsubsection{Exemple corrigé 2 - Parcours des valeurs}

\begin{python}
for valeur in exif.values():
    print(valeur)
# Affiche : Nikon D3500, 2025-06-15, 400
\end{python}

\subsubsection{Exemple corrigé 3 - Parcours des couples (clé, valeur)}

\begin{python}
for cle, valeur in exif.items():
    print(f"{cle} : {valeur}")
# Affiche :
# appareil : Nikon D3500
# date : 2025-06-15
# iso : 400
\end{python}

\subsubsection{Exemple corrigé 4 - Recherche d’une valeur}

\begin{python}
def trouver_cle(d, valeur_cherchee):
    """Renvoie la première clé associée à la valeur cherchée, ou None."""
    for cle, valeur in d.items():
        if valeur == valeur_cherchee:
            return cle
    return None
\end{python}

\begin{itemize}
\tightlist
\item
  Donnée étudiée : \texttt{exif}, recherche de \texttt{400}.
\item
  Résultat obtenu : \texttt{"iso"}.
\end{itemize}

\subsubsection{Cas limites}

\begin{itemize}
\tightlist
\item
  Itération sur un dictionnaire vide : la boucle ne s'exécute pas.
\item
  Modification du dictionnaire pendant l'itération : provoque une erreur
  \texttt{RuntimeError}. Il faut itérer sur une copie des clés si on
  veut modifier le dictionnaire.
\end{itemize}

\begin{center}\rule{0.5\linewidth}{0.5pt}\end{center}

\subsection{6. P-DATA-CONSTR-03C : Utiliser les méthodes keys, values et items}

\subsubsection{Définition}

Les méthodes \texttt{keys()}, \texttt{values()} et \texttt{items()} d'un
dictionnaire renvoient des objets vue (\emph{view objects}) qui
reflètent dynamiquement le contenu du dictionnaire. Contrairement à une
liste, une vue n'est pas une copie : elle se met à jour automatiquement
si le dictionnaire est modifié en dehors d'une itération.

\subsubsection{Formalisation}

\begin{longtable}[]{@{}lll@{}}
\toprule\noalign{}
Méthode & Type renvoyé & Contenu \\
\midrule\noalign{}
\endhead
\bottomrule\noalign{}
\endlastfoot
\texttt{d.keys()} & \texttt{dict\_keys} & toutes les clés \\
\texttt{d.values()} & \texttt{dict\_values} & toutes les valeurs \\
\texttt{d.items()} & \texttt{dict\_items} & tous les couples
\texttt{(clé,\ valeur)} \\
\end{longtable}

Ces objets sont itérables et supportent le test d'appartenance avec
\texttt{in}, mais ne sont \textbf{pas indexables} :
\texttt{d.keys(){[}0{]}} provoque une \texttt{TypeError}. Pour obtenir
une liste, on écrit \texttt{list(d.keys())}.

\subsubsection{Exemple corrigé 1 — Extraire les clés d’un dictionnaire}

\begin{python}
notes = {"Alice": 16, "Bob": 12, "Clara": 18}
cles = notes.keys()
print(cles)          # dict_keys(['Alice', 'Bob', 'Clara'])
print("Alice" in cles)  # True
\end{python}

\begin{itemize}
\tightlist
\item
  Donnée étudiée : \texttt{notes}, trois entrées.
\item
  Résultat obtenu : un objet \texttt{dict\_keys} contenant les trois
  prénoms.
\end{itemize}

\subsubsection{Exemple corrigé 2 — Calculer une moyenne avec values}

\begin{python}
notes = {"Alice": 16, "Bob": 12, "Clara": 18}
total = sum(notes.values())
moyenne = total / len(notes)
print(moyenne)  # 15.333333333333334
\end{python}

\begin{itemize}
\tightlist
\item
  Donnée étudiée : \texttt{notes.values()} renvoie
  \texttt{dict\_values({[}16,\ 12,\ 18{]})}.
\item
  Résultat obtenu : la somme vaut 46, la moyenne vaut environ 15.33.
\end{itemize}

\subsubsection{Exemple corrigé 3 — Trouver le maximum avec items}

\begin{python}
notes = {"Alice": 16, "Bob": 12, "Clara": 18}
meilleur_nom = ""
meilleure_note = -1
for nom, note in notes.items():
    if note > meilleure_note:
        meilleure_note = note
        meilleur_nom = nom
print(f"{meilleur_nom} : {meilleure_note}")  # Clara : 18
\end{python}

\begin{itemize}
\tightlist
\item
  Donnée étudiée : les couples \texttt{(nom,\ note)} via
  \texttt{items()}.
\item
  Résultat obtenu : Clara avec 18.
\end{itemize}

\subsubsection{Exemple corrigé 4 — Vue dynamique}

\begin{python}
d = {"x": 1, "y": 2}
vue = d.keys()
print(list(vue))  # ['x', 'y']
d["z"] = 3
print(list(vue))  # ['x', 'y', 'z']  — la vue reflète l'ajout
\end{python}

\begin{itemize}
\tightlist
\item
  Point clé : \texttt{vue} n'est pas une copie figée ; elle suit le
  dictionnaire.
\end{itemize}

\subsubsection{Cas limites}

\begin{itemize}
\tightlist
\item
  Sur un dictionnaire vide, \texttt{d.keys()}, \texttt{d.values()} et
  \texttt{d.items()} renvoient des vues vides (itérables, longueur 0).
\item
  Modifier le dictionnaire \textbf{pendant} une itération sur une vue
  provoque \texttt{RuntimeError}. Solution : itérer sur
  \texttt{list(d.keys())} si des modifications sont nécessaires.
\end{itemize}

\begin{center}\rule{0.5\linewidth}{0.5pt}\end{center}

\subsection{Synthèse}

\begin{longtable}[]{@{}
  >{\raggedright\arraybackslash}p{(\columnwidth - 4\tabcolsep) * \real{0.2778}}
  >{\raggedright\arraybackslash}p{(\columnwidth - 4\tabcolsep) * \real{0.3056}}
  >{\raggedright\arraybackslash}p{(\columnwidth - 4\tabcolsep) * \real{0.4167}}@{}}
\toprule\noalign{}
\begin{minipage}[b]{\linewidth}\raggedright
Capacité
\end{minipage} & \begin{minipage}[b]{\linewidth}\raggedright
Structure
\end{minipage} & \begin{minipage}[b]{\linewidth}\raggedright
Opération clé
\end{minipage} \\
\midrule\noalign{}
\endhead
\bottomrule\noalign{}
\endlastfoot
P-DATA-CONSTR-01 & Tuple & \texttt{return\ a,\ b} \\
P-DATA-CONSTR-02B & Liste &
\texttt{{[}expr\ for\ var\ in\ iterable{]}} \\
P-DATA-CONSTR-02C & Liste de listes & \texttt{matrice{[}i{]}{[}j{]}} \\
P-DATA-CONSTR-03A & Dictionnaire & \texttt{d{[}cle{]}\ =\ valeur} \\
P-DATA-CONSTR-03B & Dictionnaire & \texttt{for\ k,\ v\ in\ d.items()} \\
P-DATA-CONSTR-03C & Dictionnaire & \texttt{d.keys()},
\texttt{d.values()}, \texttt{d.items()} \\
\end{longtable}

\subsection{Situation-problème}

Un programme de gestion de notes doit stocker, pour chaque élève, son
nom et ses notes dans différentes matières. Quelle structure de données
choisir et comment organiser les accès ?

\subsection{Activité d’entrée}

Écrire une fonction Python qui reçoit deux nombres et renvoie à la fois
leur somme et leur produit. Tester avec print(resultat{[}0{]},
resultat{[}1{]}).

\subsection{Exercices}

Exercices de manipulation de tuples, listes en compréhension, matrices
et dictionnaires.

\subsection{Corrigé}

Les corrigés détaillés sont dans
P04\_corrige\_types\_construits\_complement.md.

\subsection{Erreurs fréquentes}

\begin{itemize}
\tightlist
\item
  EF1 : confondre tuple (immutable) et liste (mutable) lors du retour de
  fonction.
\item
  EF2 : oublier les crochets dans une compréhension de liste.
\item
  EF3 : accéder à une clé inexistante d'un dictionnaire sans utiliser
  get().
\end{itemize}

\subsection{Remédiation}

Exercice de remédiation : écrire une fonction min\_max qui renvoie le
minimum et le maximum d'une liste sous forme de tuple.

\subsection{Différenciation}

\begin{itemize}
\tightlist
\item
  Socle : exercices de base.
\item
  Standard : exercices complets.
\item
  Expert : exercice d'approfondissement.
\end{itemize}

\subsection{Critères de réussite}

Fonction renvoyant un tuple vérifiée. Compréhension de liste correcte.
Matrice accessible par double index. Dictionnaire itéré avec
keys/values/items.

Critère de validation : chaque réponse est vérifiable par un contrôle
explicite.
